\clearpage

\section*{Acknowledgement}
S.P. and A.M. are supported in part by the Engineering and Physical Sciences Research Council (EPSRC) under Grant EP/X028909/1 and Oxford Internet Institute’s Research Programme funded by the Dieter Schwarz Foundation. R.O.K. is supported by the Clarendon Scholarship and the Jesus College Old Members’ Scholarship. E.S. acknowledges support in part from the AI2050 program at Schmidt Sciences (Grant No. G-22-64476). E.S. and A.C. acknowledge that this study is funded by the National Institute for Health Research (NIHR) Health Protection Research Unit in Health Analytics \& Modelling (NIHR207404), a partnership between UK Health Security Agency (UKHSA), London School of Hygiene \& Tropical Medicine, and Imperial College of Science, Technology, \& Medicine. The views expressed are those of the author(s) and not necessarily those of the NIHR, UKHSA, or the Department of Health and Social Care. The authors thank Saverio Trioni for helpful conversations. The authors would like to thank the Pathogen Epidemiology Review Group (PERG), School of Public Health, Imperial College London, for their support and eagerness to contribute to this project.

% \section*{Impact Statement}

% This paper presents work whose goal is to advance the field of Machine Learning. There are many potential societal consequences of our work, particularly in the context of public health and scientific evidence synthesis. By demonstrating the feasibility of automating substantial components of systematic literature reviews in epidemiology, this work has the potential to significantly reduce the time and cost required to synthesise scientific evidence. Such efficiency gains may enable faster translation of research findings into policy and practice, support the maintenance of ``living" reviews, and improve preparedness and response during public health emergencies.

% Automated evidence synthesis systems are not without risks. Errors, biases, or omissions introduced during retrieval, screening, or extraction could propagate into downstream analyses and decision-making if not carefully managed. To mitigate these risks, our system is explicitly designed for human-in-the-loop use, prioritising high recall, structured outputs, and transparent validation over full autonomy. We do not advocate replacing expert judgment, but rather augmenting expert workflows to reduce manual burden while preserving accountability.
% Our emphasis on open-weight models, schema-constrained extraction, and explicit evaluation is intended to support responsible development and scrutiny. We believe that, when carefully designed and supported by human-in-the-loop, such tools can meaningfully enhance scientific practice while maintaining the rigour and trust essential to public health research.